\section{Технический проект}
\subsection{Выбор технологии проектирования}

Выбор технологий обусловлен целями дипломного проекта в рамках темы «ПИС регистрации данных пользователей»: наглядная клиент-серверная архитектура, явная регистрация и учёт пользователей и связанных сущностей, ролевая модель, регламентированный workflow операций (заявок) и безопасность данных при ограниченных сроках разработки.

\textbf{Тип системы:} веб-приложение (SPA), клиент-серверное взаимодействие через REST API Supabase. Данные и аутентификация размещаются в облаке Supabase; приложение не содержит собственного серверного кода, что упрощает развёртывание и демонстрацию.

\textbf{Язык и фронтенд:} TypeScript и Vue 3 (Composition API) обеспечивают типизацию, предсказуемую структуру компонентов и переиспользуемую логику (composables). Сборка через Vite даёт быструю разработку и оптимизированную production-сборку. Разделение на домены (auth, requests, customers, vehicles, users, audit) позволяет чётко выделить зоны ответственности за регистрацию данных пользователей и связанных сущностей.

\textbf{Хранение данных:} реляционная СУБД PostgreSQL (Supabase). Структурированные сущности (пользователи, профили с ролями, клиенты, автомобили, заявки, аудит) хорошо отображаются в таблицах и связях; RLS реализуется средствами PostgreSQL и обеспечивает разграничение доступа на уровне строк без дублирования логики в приложении.

\textbf{Безопасность:} аутентификация и сессии — Supabase Auth; разграничение доступа — RLS-политики в БД плюс проверки в приложении (guards при навигации, скрытие UI по ролям). Многоуровневая защита гарантирует, что даже при обходе интерфейса пользователь не получит доступ к чужим данным.

\textbf{Масштабируемость и поддерживаемость:} модульная доменная структура позволяет добавлять новые домены (например, отчёты, уведомления) и функции без ломки существующих модулей. Зависимости направлены внутрь: UI → composables/store → API → Supabase, что соответствует принципам чистой архитектуры и облегчает тестирование и сопровождение.

Критерии выбора: возможность реализовать полный цикл регистрации и учёта данных пользователей (от входа до аудита), использование открытых и широко применяемых технологий, пригодность для защиты диплома и последующего расширения или переноса примера на другую отрасль.

\subsection{Выбор средств разработки}

\textbf{Vue 3 (Composition API)} — фреймворк для построения SPA. Преимущества: реактивность, компонентный подход, богатая экосистема (Router, Pinia), удобство для типизации с TypeScript. В проекте используется только Composition API и \textquotedbl <script setup>\textquotedbl\ , что обеспечивает единообразие кода и хорошую читаемость. Компоненты отвечают за презентацию; бизнес-логика вынесена в composables и store.

\textbf{TypeScript:} — статическая типизация снижает ошибки на этапе разработки и улучшает поддержку рефакторинга. Типы доменов (заявка, клиент, автомобиль, пользователь, запись аудита) описаны в файлах \textquotedbl types.ts\textquotedbl\ внутри каждого домена, что делает контракты данных явными и устойчивыми к изменениям.

\textbf{Pinia:} — хранилище состояния для Vue 3. Используется для доменного состояния (списки заявок, клиентов, выбранная заявка и т.д.); UI-состояние (модальные окна, вкладки, флаги загрузки) остаётся в компонентах. Такое разделение упрощает отладку и не раздувает глобальное состояние.

\textbf{Vue Router:} — маршрутизация. Lazy loading маршрутов через динамический \textquotedbl import()\textquotedbl\ уменьшает размер начальной загрузки. Глобальный guard проверяет аутентификацию; для защищённых маршрутов используется role guard (массив допустимых ролей), что напрямую поддерживает требования регистрации данных пользователей и разграничения доступа.

\textbf{Supabase:} — Backend-as-a-Service: PostgreSQL, аутентификация, REST API, генерация типизированного клиента для JavaScript/TypeScript. RLS реализуется на стороне БД; триггеры позволяют автоматически заполнять журнал аудита при изменениях в ключевых таблицах. Преимущества для дипломного проекта: быстрый старт, не требуется развёртывание отдельного бэкенд-сервера, наглядность политик RLS и схемы данных.

\textbf{Vite:} — сборщик и dev-сервер. Обеспечивает быстрый холодный старт и горячую перезагрузку при разработке, а также оптимизированную сборку для production с разбиением кода по маршрутам.

\textbf{PrimeVue:} — набор UI-компонентов (таблицы, формы, кнопки, диалоги, сообщения). Ускоряет разработку интерфейса при сохранении единообразного вида и доступности.

Стек Vue 3 + TypeScript + Pinia + Vue Router + Supabase + Vite в совокупности обеспечивает единообразие кода, разделение слоёв (UI, логика, данные), соответствие правилам проекта (единый API-слой домена, доменная структура, RLS) и пригодность системы для демонстрации механизмов регистрации и учёта данных пользователей.

\subsection{Компоненты приложения}

(API), состояние (store), бизнес-логику (composables), компоненты и страницы (views). Общие элементы вынесены в \textquotedbl lib\textquotedbl\ (константы, клиент Supabase, утилиты), \textquotedbl router\textquotedbl\ (маршруты и guards), \textquotedbl layouts\textquotedbl\, \textquotedbl views\textquotedbl\ (общие страницы — Dashboard, NotFound). Такая организация соответствует теме ПИС регистрации данных пользователей: домен auth отвечает за учётные записи и роли, домены customers, vehicles, requests — за регистрацию и учёт связанных сущностей, домен users — за управление пользователями системы (для администратора), домен audit — за отображение журнала изменений.

На рисунке \ref{component-diagram:image} представлена диаграмма компонентов.

\begin{figure}[H]
	\includegraphics[width=1\linewidth]{component-diagram}
	\caption{Диаграмма компонентов приложения}
	\label{component-diagram:image}
\end{figure}

\textbf{Описание компонентов и слоёв.}

\begin{itemize}
	\item \textbf{Vue Router — маршруты:} "</login">, "</dashboard">, "</requests">, "</requests/:id">, "</customers">, "</vehicles">, "</users">, "</audit">. Глобальный guard проверяет аутентификацию; для защищённых маршрутов используется role guard (массив допустимых ролей), что ограничивает доступ к разделам в соответствии с регистрацией данных пользователей и их ролями;
	\item \textbf{Домен Auth:} страница входа (LoginPage), store сессии и профиля (в том числе роль текущего пользователя), composable useAuth, guards (createRoleGuard). Взаимодействует с Supabase Auth и таблицей профилей пользователей;
	\item \textbf{Домен Requests:} страницы списка заявок и карточки заявки, форма создания/редактирования, store заявок, API requests-api (и при необходимости request-comments-api), composables useRequestWorkflow, useRequestStatusPermissions. Бизнес-логика переходов статусов и проверка прав вынесены в composables и модуль workflow; API — единственная точка обращения к Supabase по заявкам, что обеспечивает предсказуемость и тестируемость;
	\item \textbf{Домен Customers:} список клиентов, форма клиента, store, customers-api. Аналогично — домен Vehicles (список автомобилей, форма, store, vehicles-api);
	\item \textbf{Домен Users:} страница пользователей системы (доступна только администратору), назначение ролей, users-api. Прямо реализует требование регистрации и учёта данных пользователей со стороны администратора;
	\item \textbf{Домен Audit:} страница журнала изменений, фильтры по сущности, действию, дате, audit-api. Отображает зафиксированные системой изменения данных;
	\item \textbf{Pinia stores:} хранят доменное состояние (списки заявок, клиентов, выбранная заявка и т.д.); загрузка данных выполняется через вызовы API из store или из компонентов. UI-состояние в store не дублируется;
	\item \textbf{Supabase:} аутентификация (Auth), таблицы PostgreSQL (user\_profiles, customers, vehicles, service\_requests, audit\_logs и др.), RLS-политики, при необходимости RPC для сложных обновлений (смена статуса заявки, назначение мастера, «взять в работу»), триггеры для записи в audit\_logs
\end{itemize}

Процесс работы: пользователь входит → Router с учётом guards направляет на Dashboard или разделы по роли → компоненты запрашивают данные через store/API → API обращается к Supabase → RLS фильтрует строки по текущему пользователю и роли → ответ отображается в UI. Смена статуса заявки: пользователь нажимает действие → composable проверяет допустимость перехода по графу и права пользователя → store вызывает API смены статуса → Supabase (RPC или update) обновляет данные, триггер при необходимости пишет в audit\_logs → обновлённые данные отображаются в интерфейсе. Таким образом достигаются разделение ответственности по слоям, масштабируемость за счёт доменов, тестируемость (API и composables можно тестировать изолированно), соответствие принципам SOLID и правилам проекта.

\subsubsection{Описание работы с внешним API (Supabase)}

Приложение использует Supabase как внешний сервис: аутентификация, БД и REST API.

\textbf{Инициализация клиента Supabase}

Клиент создаётся один раз с URL проекта и anon-ключом; все запросы к БД и Auth идут через него. Переменные окружения: \textquotedbl VITE\_SUPABASE\_URL\textquotedbl\, \textquotedbl VITE\_SUPABASE\_ANON\_KEY\textquotedbl\. Типизация через \textquotedbl Database\textquotedbl\ обеспечивает соответствие ответов API схеме БД.

\begin{figure}[H]
	\includegraphics[width=1\linewidth]{init-client-supabase}
	\caption{Инициализация клиента Supabase}
	\label{init-client-supabase:image}
\end{figure}

\textbf{Запросы к данным (заявки)}

Весь доступ к таблице заявок инкапсулирован в API домена requests. Пример: получение списка заявок с развёрнутыми связями (клиент, автомобиль, менеджер, мастер) и изменение статуса через RPC. RPC обеспечивает выполнение бизнес-правил на стороне БД (проверка допустимости перехода, запись в аудит).

\begin{figure}[H]
	\includegraphics[width=1\linewidth]{zayavki}
	\caption{Запросы к данным (заявки)}
	\label{zayavki:image}
\end{figure}

\textbf{Workflow и права (composable)}

Бизнес-логика допустимых переходов статусов и проверка прав по роли вынесены в composable. Константы переходов заданы в \textquotedbl src/lib/constants.ts\textquotedbl\ (REQUEST\_STATUS\_TRANSITIONS).

Ниже — фрагмент composable для смены статуса. 

\begin{figure}[H]
	\includegraphics[width=1\linewidth]{workflow-permissions-composable}
	\caption{Workflow и права (composable)}
	\label{workflow-permissions-composable:image}
\end{figure}

\textbf{Константы переходов статусов}

\begin{figure}[H]
	\includegraphics[width=1\linewidth]{status-transition-const}
	\caption{Константы переходов статусов}
	\label{status-transition-const:image}
\end{figure}

\newpage
\textbf{Обработка ошибок}

\begin{itemize}
	\item Ошибки сети и ответы Supabase (поле error) обрабатываются в API: выбрасывается исключение с сообщением; вызывающий код (store или компонент) отображает сообщение пользователю и при необходимости сбрасывает состояние загрузки;
	\item При недопустимом переходе статуса или отсутствии прав composable выбрасывает ошибку; UI заранее блокирует недоступные кнопки на основе тех же правил (canUserChangeStatus, getAvailableTransitions), что обеспечивает согласованность логики и интерфейса.
\end{itemize}

Использование единого API-слоя и явного workflow обеспечивает предсказуемость обмена с Supabase и удобство сопровождения кода в рамках ПИС регистрации данных пользователей.

\subsection{Проектирование потоков данных и состояния}

В программно-информационной системе регистрации данных пользователей важен не только состав компонентов, но и то, как данные проходят через приложение. Пользователь видит формы, таблицы и кнопки, однако за каждым действием находится цепочка: пользовательский интерфейс передаёт событие в бизнес-логику, бизнес-логика обращается к доменному состоянию, store вызывает API-слой, API-слой работает с Supabase, а PostgreSQL применяет политики RLS. Такой подход делает поведение системы предсказуемым и позволяет объяснить, где именно выполняется каждая ответственность.

Ключевое решение технического проекта — <<не обращаться к Supabase напрямую из UI-компонентов>>. Компонент не должен знать, из какой таблицы берутся данные, какие RPC-функции вызываются и какие детали есть у внешнего API. Компонент отвечает за отображение, формы и пользовательские события. Бизнес-логика находится в composables, доменное состояние — в Pinia, доступ к данным — в API-слое домена.

На рисунке \ref{data-flow-app:image} представлена диаграмма потока данных в приложении.

\begin{figure}[H]
	\centering
	\includegraphics[
	height=0.9\textheight,
	keepaspectratio
	]{data-flow-app}
	\caption{Диаграмма потока данных в приложении}
	\label{data-flow-app:image}
\end{figure}

\subsubsection{Маршрут пользовательского действия}

Рассмотрим типовой маршрут действия на примере смены статуса заявки. Пользователь открывает карточку заявки и нажимает кнопку, например <<Начать работу>> или <<Завершить>>. Компонент страницы не выполняет бизнес-правила самостоятельно, а вызывает функцию из composable workflow. Composable получает текущую заявку из store, проверяет допустимость перехода по графу статусов и права текущего пользователя. Если переход допустим, store вызывает метод API-слоя, а API-слой обращается к RPC-функции Supabase.

Такой маршрут действия позволяет разделить ответственность:

\begin{itemize}
	\item UI-компонент отвечает за отображение кнопки, диалога и результата действия;
	\item Composable отвечает за бизнес-правила и координацию операции;
	\item Store отвечает за доменное состояние и обновление локального представления данныхStore отвечает за доменное состояние и обновление локального представления данных;
	\item API-слой отвечает за обмен с Supabase и обработку ошибок внешнего сервиса;
	\item База данных отвечает за хранение, RLS и аудит;
\end{itemize}

Такой подход особенно важен для темы <<Программно-информационная система регистрации данных пользователей>>, потому что каждое действие должно быть связано с конкретным пользователем, ролью и правами доступа. Если пользователь не имеет права выполнить действие, операция блокируется до изменения данных.

\subsubsection{Управление серверным и локальным состоянием}

В приложении различаются два вида состояния: серверное и локальное. Серверное состояние — это данные, которые являются частью предметной области и должны быть сохранены в БД: пользователи, клиенты, автомобили, заявки, комментарии и записи аудита. Локальное состояние — это временное состояние интерфейса: открыт ли диалог, выбран ли фильтр, идёт ли отправка формы, отображается ли сообщение об ошибке.

Серверное состояние хранится в доменных Pinia stores. Например, store заявок содержит список заявок, выбранную заявку, флаги загрузки и методы для загрузки или изменения данных. При этом store не содержит детали модальных окон и формы редактирования: такое состояние остаётся внутри Vue-компонентов. Это снижает связанность и предотвращает превращение глобального состояния в <<god-store>>.

Производные данные не дублируются, а вычисляются через <<computed>>. Например, активные заявки, доступные действия по статусу или видимость кнопок можно получать из текущих данных, роли пользователя и правил workflow. Это уменьшает риск рассинхронизации: если изменился статус заявки, интерфейс автоматически пересчитывает доступные действия.

\subsubsection{Обработка ошибок и состояний загрузки}

Асинхронные операции в системе выполняются с учётом возможных ошибок: сетевой сбой, ошибка Supabase, отсутствие прав, недопустимый переход статуса, конфликт при назначении мастера. Ошибка не должна оставлять интерфейс в неопределённом состоянии. Поэтому для операций загрузки используются флаги <<isLoading>>, для действий — отдельные флаги блокировки кнопок, а для сообщений пользователю — понятные текстовые уведомления.

При загрузке списка или карточки заявок интерфейс показывает состояние ожидания. Если запрос завершился ошибкой, пользователю отображается сообщение, а технические детали не раскрываются. При успешном выполнении операции store обновляет данные локально или повторно загружает актуальную запись из Supabase. Для критичных операций, таких как смена статуса или взятие заявки в работу, предпочтительно использовать RPC-функции, потому что они позволяют выполнять проверку и изменение данных атомарно на стороне БД.

Такой подход повышает надёжность интерфейса: пользователь видит, что действие выполняется, не может случайно отправить одно и то же действие несколько раз и получает понятную обратную связь при ошибке.

\subsection{Надёжность, производительность и масштабирование}

Технический проект должен учитывать не только текущий объём данных, но и возможность роста системы. Даже если дипломный проект демонстрируется на ограниченном наборе записей, архитектура должна оставаться пригодной для увеличения числа пользователей, заявок, клиентов и связанных сущностей. Поэтому в проект заложены решения, влияющие на надёжность загрузки, производительность интерфейса и потенциал масштабирования.

\subsubsection{Надёжность загрузки и восстановления данных}

Надёжность загрузки данных обеспечивается несколькими уровнями. На уровне приложения запросы выполняются через единый API-слой, поэтому обработка ошибок и формат возвращаемых данных единообразны. На уровне store фиксируются состояния загрузки и ошибки. На уровне БД RLS гарантирует, что пользователь получает только допустимые данные, а не весь массив записей с последующей фильтрацией на клиенте.

Если пользователь обновляет страницу или повторно входит в систему, приложение заново инициализирует сессию, загружает профиль и роль, после чего маршруты и данные становятся доступны в соответствии с правами. Это важно для сценариев, когда сессия уже существует, но приложение открыто в новой вкладке. Такой механизм не зависит от локально сохранённого состояния и восстанавливает приложение из серверного источника правды.

Для операций изменения данных важно, что после успешной мутации приложение может повторно запросить актуальную запись. Например, при взятии заявки в работу API возвращает обновлённую заявку или инициирует повторную загрузку по ID. Это защищает интерфейс от устаревших данных и конфликтов, если другая роль уже изменила запись.

\subsubsection{Оптимизация списков, поиска и маршрутов}

В системе используются списки заявок, клиентов, автомобилей и пользователей. Такие экраны потенциально могут содержать большое количество записей, поэтому проект предусматривает фильтрацию, сортировку и ограничение количества загружаемых данных. Для заявок реализован API списка с параметрами поиска, статусов, сортировки, "<limit"> и "<offset">. Это позволяет не загружать весь набор заявок сразу и уменьшает нагрузку на браузер и БД.

Поиск по заявкам учитывает не только описание, но и связанные сущности — клиента и автомобиль. Такой подход удобен пользователю, потому что заявка может быть найдена по ФИО, телефону, госномеру, марке или модели. При этом сама логика поиска находится в API-слое домена, а не в компоненте, что сохраняет архитектурную чистоту.

Маршруты приложения загружаются лениво через динамический "<import()">. Это означает, что код страницы пользователей, аудита или карточки заявки не обязательно попадает в начальный пакет приложения. Пользователь быстрее открывает основной экран, а остальные части загружаются по мере переходов. Для SPA это важный фактор производительности и удобства.

На уровне БД производительности способствуют индексы по полям, которые участвуют в фильтрах, связях и RLS: статус заявки, идентификатор мастера, клиент, автомобиль, дата создания. Индексы позволяют быстрее выполнять выборки и политики доступа, особенно при росте числа записей.

\subsubsection{Потенциал масштабирования системы}

Система масштабируется прежде всего за счёт доменной архитектуры. Если необходимо добавить новый блок, например <<Отчёты>>, <<Уведомления>>, <<Склад>> или <<Платежи>>, можно создать отдельный домен с собственными типами, API-слоем, store, composables и страницами. При этом существующие домены не требуется переписывать, так как взаимодействие между ними идёт через явные точки: типы, API и маршруты.

Supabase обеспечивает масштабирование backend-части за счёт PostgreSQL, REST API, Auth и RLS. Для дипломного проекта этого достаточно, чтобы показать полноценную клиент-серверную архитектуру без отдельного серверного приложения. В дальнейшем при необходимости можно вынести часть логики в Edge Functions или отдельный backend, но текущая архитектура уже отделяет доступ к данным через API-слой, поэтому такая миграция не потребует переписывать UI-компоненты.

Масштабирование также возможно по ролям. Сейчас реализованы admin, manager и technician. В дальнейшем можно добавить роль <<customer>> или <<auditor>>, расширить RLS-политики и маршруты, не меняя базовую схему взаимодействия. Это подтверждает, что проект реализует не только конкретный пример CRM, но и основу ПИС регистрации данных пользователей.

\subsection{Информационная безопасность и эксплуатация технического проекта}

Безопасность является одной из ключевых частей технического проекта, потому что система работает с пользовательскими данными, ролями, заявками и историей изменений. Ошибка в разграничении доступа может привести к просмотру чужих данных, изменению роли, некорректной смене статуса или удалению записей. Поэтому защита строится многоуровнево: интерфейс, маршруты, бизнес-логика, API-слой и RLS в базе данных.

\subsubsection{Профиль угроз}

Для системы выделены следующие основные угрозы:

\begin{enumerate}
	\item \textbf{Несанкционированный доступ к разделам приложения.} Пользователь без нужной роли может попытаться открыть URL административного раздела напрямую.
	
	\item \textbf{Прямое обращение к Supabase API.} Пользователь может попытаться обойти интерфейс и отправить запрос к REST API вручную.
	
	\item \textbf{Подмена роли на клиенте.} Так как frontend выполняется в браузере, нельзя доверять только клиентским проверкам.
	
	\item \textbf{Некорректная смена статуса заявки.} Пользователь может попытаться перевести заявку в недопустимый статус или изменить заявку, к которой не имеет отношения.
	
	\item \textbf{Утечка конфигурации.} Публичный anon-ключ Supabase может быть виден в клиентском приложении, поэтому он не должен давать административных прав без RLS.
	
	\item \textbf{Потеря прозрачности изменений.} Без аудита невозможно установить, кто изменил важные данные.
\end{enumerate}	

\subsubsection{Меры защиты}

Основная мера защиты --- Supabase Auth и RLS. После входа пользователь получает сессию, а база данных применяет политики доступа с учётом \texttt{auth.uid()} и роли из профиля. Это означает, что даже при прямом запросе к API пользователь увидит только те строки, которые разрешены политиками.

На уровне маршрутов используется \texttt{authGuard} и role guard. Они предотвращают переход на закрытые страницы, например страницу управления пользователями для неадминистратора. На уровне интерфейса недоступные действия скрываются или блокируются: мастер не видит административные действия, а менеджер не может выполнять операции, которые разрешены только администратору.

Критичные операции выполняются через RPC-функции Supabase. Например, смена статуса заявки или назначение мастера может быть реализовано как единая операция на стороне БД. Это снижает риск частичного обновления данных и позволяет централизованно проверять бизнес-правила. Если операция успешна, триггеры БД фиксируют изменение в \texttt{audit\_logs}, сохраняя историю действий.

Переменные окружения (\texttt{VITE\_SUPABASE\_URL}, \texttt{VITE\_SUPABASE\_ANON\_KEY}) хранятся вне исходного кода. Anon-ключ является публичным ключом клиента, поэтому безопасность не должна основываться на его сокрытии. Реальная защита обеспечивается RLS, ролями, RPC и ограничением операций на уровне БД.

\subsection{Разграничение прав доступа}

Для обеспечения информационной безопасности и разграничения зон ответственности в программно-информационной системе реализована ролевая модель доступа, включающая роли "<Администратор">, "<Менеджер"> и "<Мастер">. Разграничение прав выполнено как на уровне клиентского приложения, так и на уровне базы данных с использованием механизмов аутентификации Supabase Auth и политик безопасности строк (Row Level Security, RLS).

Администратор обладает полным набором прав на все сущности системы. Пользователь данной роли может создавать, редактировать и удалять записи пользователей, клиентов, автомобилей и заявок, назначать роли пользователям, просматривать журнал изменений, а также выполнять любые действия, предусмотренные функциональностью системы.

Менеджер имеет доступ к операциям, связанным с обработкой заявок и ведением клиентской базы. Пользователь данной роли может создавать и редактировать заявки, регистрировать новых клиентов и автомобили, назначать исполнителей, изменять статусы заявок в пределах установленных бизнес-правил, а также просматривать журнал изменений. При этом менеджер не имеет доступа к управлению учётными записями пользователей и назначению ролей.

Мастер (техник) обладает ограниченным набором прав, необходимым для выполнения производственных задач. Пользователь данной роли может просматривать назначенные ему заявки, принимать их в работу и изменять статус в соответствии с регламентированным жизненным циклом. Доступ к управлению клиентами, пользователями и системными настройками для данной роли запрещён.

При входе в систему выполняется аутентификация пользователя посредством Supabase Auth. После успешного входа определяется роль пользователя, на основании которой формируется пользовательский интерфейс: отображаются только разрешённые разделы, команды и элементы управления. Дополнительно политики RLS ограничивают доступ к данным на уровне базы данных, исключая возможность получения или изменения информации в обход интерфейса приложения.

\subsection{Обеспечение надёжности и целостности данных}

Требования к надёжности и целостности данных, сформулированные в техническом задании, реализованы в системе с использованием механизмов реляционной базы данных PostgreSQL и средств платформы Supabase.

Целостность данных обеспечивается за счёт использования первичных и внешних ключей, ограничений целостности и проверки корректности вводимых данных. Между сущностями системы установлены логические связи: каждый автомобиль привязан к определённому клиенту, каждая заявка связана с клиентом, автомобилем и ответственными пользователями. Попытки создания записей, нарушающих установленные связи, блокируются средствами СУБД.

Корректность вводимых данных контролируется как на стороне клиентского приложения, так и на стороне базы данных. В интерфейсе реализована проверка обязательности заполнения полей, форматов телефонных номеров и адресов электронной почты, а также допустимости выполняемых операций. Дополнительная проверка осуществляется при выполнении запросов к базе данных.

Особое внимание уделено обеспечению корректности жизненного цикла заявок. Изменение статусов осуществляется в соответствии с заранее определённым набором допустимых переходов. Попытки выполнения недопустимых операций отклоняются системой независимо от источника запроса.

Для повышения надёжности работы приложения предусмотрена обработка исключительных ситуаций, возникающих при взаимодействии с сервером и внешними сервисами. Пользователю отображаются информативные сообщения об ошибках без раскрытия внутренней структуры системы и деталей реализации.

Дополнительным механизмом обеспечения достоверности данных является журнал изменений (аудит), в котором автоматически фиксируются сведения о действиях пользователей над ключевыми сущностями системы. Для каждой операции сохраняются тип изменённого объекта, выполненное действие, идентификатор пользователя и время выполнения операции. Наличие журнала аудита позволяет отслеживать историю изменений, проводить анализ спорных ситуаций и повышает уровень ответственности пользователей при работе с системой.

Использование перечисленных механизмов обеспечивает устойчивую работу программно-информационной системы, сохранность и непротиворечивость данных, а также соответствие требованиям информационной безопасности и надёжности, предъявляемым к современным веб-приложениям.

\subsubsection{Эксплуатация и развертывание}

Эксплуатация приложения включает локальный запуск, настройку Supabase, сборку production-версии и деплой статического SPA. Для локальной разработки необходимо установить зависимости, создать файл \texttt{.env}, указать \texttt{VITE\_SUPABASE\_URL} и \texttt{VITE\_SUPABASE\_ANON\_KEY}, выполнить SQL-скрипты Supabase и запустить dev-сервер.

\begin{verbatim}
	npm install
	npm run dev
	npm run build
	npm run typecheck
\end{verbatim}

При развертывании на Vercel проект собирается как Vite SPA. Переменные окружения задаются в панели Vercel, после чего при каждом push в основную ветку Vercel запускает сборку и публикует новый билд. После деплоя необходимо проверить главную страницу, страницу входа, защищённые маршруты (\texttt{/dashboard}, \texttt{/requests}, \texttt{/customers}, \texttt{/vehicles}, \texttt{/users}, \texttt{/audit}) и доступ к Supabase.

Для эксплуатации важны следующие проверки: наличие переменных окружения, выполненные SQL-скрипты схемы и RLS, созданные тестовые пользователи с ролями, доступность Supabase, корректность маршрутов после обновления страницы. Так как приложение является SPA, хостинг должен перенаправлять неизвестные маршруты на \texttt{index.html}, иначе прямое открытие \texttt{/requests} или \texttt{/dashboard} может возвращать ошибку 404.

Таким образом, раздел эксплуатации показывает, что проект готов не только к локальному запуску, но и к размещению в production-среде с понятным набором проверок и зависимостей.